% Chapter 5. Estimation and prediction results from the frozen R=5,000 archive.
\chapter{Estimation and Prediction Results}
\label{chap:estimation-results}

The fixed experiment from Chapter~\ref{chap:simulation-design} is analysed
without post hoc changes to its tuning rules or diagnostics.  The results first
separate slope recovery from independent prediction and then use selected
dimensions, tail behaviour, and paired raw--Arnoldi fits to distinguish
statistical regularisation from numerical stabilisation.  The central issue is
not a single method ranking, but whether the target geometry and numerical
representation favour the same amount and form of regularisation.

\section{Archive verification and reporting conventions}
\label{sec:results-verification}

All $60{,}000$ fits enter the primary ISE, MSPE, component-count, and tail
summaries.  Function-valued replay excludes only the $445$ raw fits flagged by
the fixed numerical-stability rule, never on the basis of their observed loss;
the adverse values therefore remain part of the reported performance of the
complete raw procedure.

Means, medians, and upper quantiles answer different questions in this
experiment.  The mean estimates the risk of using the complete procedure,
including rare adverse outcomes.  The median describes a typical
replication, while the $90$th and $99$th percentiles describe progressively
less common losses without allowing a single catastrophic value to dominate.
For the numerically stable methods, the Monte Carlo standard errors of mean
ISE range from $0.0018$ to $0.0475$, and those of mean MSPE range from
$0.0024$ to $0.0088$.  The exceptionally large standard errors attached to
some raw-FPLS means are themselves evidence of a heavy numerical tail.

\section{Overall estimation and prediction accuracy}
\label{sec:results-overall-accuracy}

Table~\ref{tab:results-primary-risk} gives the primary full-archive results.
All $5{,}000$ replications enter each row, including every numerically flagged
raw fit.  The boxplots in Figure~\ref{fig:results-error-boxplots} display the
same observations on logarithmic axes, making both central performance and
the complete upper tails visible.

\begin{table}[htbp]
    \centering
    \scriptsize
    \caption[Full-archive slope and prediction errors]{Full-archive slope-estimation and prediction results.  Means,
    medians, and $90$th percentiles are calculated over all $5{,}000$
    replications for each model--method pair.}
    \label{tab:results-primary-risk}
    \begin{tabular}{@{}llrrrrrr@{}}
        \toprule
        & & \multicolumn{3}{c}{Slope ISE}
          & \multicolumn{3}{c}{Test MSPE} \\
        \cmidrule(lr){3-5}\cmidrule(l){6-8}
        Model & Method & Mean & Median & $p_{90}$ & Mean & Median & $p_{90}$ \\
        \midrule
        Model 1 & CG--FPLS       & 0.633 & 0.597 & 0.910 & 1.270 & 1.230 & 1.617 \\
                & Raw FPLS       & 334.298 & 0.539 & 0.990 & 21.224 & 1.114 & 1.357 \\
                & Arnoldi FPLS   & 0.760 & 0.539 & 0.987 & 1.130 & 1.113 & 1.356 \\
                & FPCR           & 0.670 & 0.532 & 0.952 & 1.126 & 1.112 & 1.357 \\
        \addlinespace
        Model 2 & CG--FPLS       & 5.506 & 5.272 & 7.354 & 1.999 & 1.861 & 2.753 \\
                & Raw FPLS       & 7.594 & 5.384 & 9.949 & 1.518 & 1.428 & 1.793 \\
                & Arnoldi FPLS   & 6.065 & 5.345 & 8.736 & 1.442 & 1.419 & 1.763 \\
                & FPCR           & 5.662 & 5.365 & 7.392 & 1.495 & 1.472 & 1.836 \\
        \addlinespace
        Model 3 & CG--FPLS       & 0.447 & 0.426 & 0.608 & 1.221 & 1.176 & 1.534 \\
                & Raw FPLS       & 1.219 & 0.446 & 0.861 & 1.154 & 1.114 & 1.365 \\
                & Arnoldi FPLS   & 0.670 & 0.446 & 0.856 & 1.131 & 1.114 & 1.363 \\
                & FPCR           & 0.615 & 0.463 & 0.930 & 1.127 & 1.113 & 1.349 \\
        \bottomrule
    \end{tabular}
\end{table}

\begin{figure}[H]
    \centering
    \includegraphics[width=\textwidth]{babii_style_error_boxplots_with_tails.pdf}
    \caption[Full empirical error distributions]{Complete empirical distributions of slope ISE (left column) and
    independent test MSPE (right column), with one row per model.  Boxes use the conventional $1.5$
    interquartile-range whiskers, all displayed axes are logarithmic, and red
    open diamonds identify raw-FPLS fits satisfying the numerical-instability
    rule.  The labels ``CG'', ``Raw'', and ``Arnoldi'' denote the released-code
    CG--FPLS, raw FPLS, and Arnoldi FPLS procedures.  No tail observation is removed.}
    \label{fig:results-error-boxplots}
\end{figure}

\subsection{Model 1: rapidly concentrated predictive signal}

In Model~1, the typical slope losses are close.  FPCR has the smallest median
ISE at $0.532$, followed by raw and Arnoldi FPLS at approximately $0.539$;
the CG--FPLS median is $0.597$.  The ranking of the stable methods changes
when the arithmetic mean is used.  CG--FPLS has mean ISE $0.633$, compared
with $0.670$ for FPCR and $0.760$ for Arnoldi FPLS, because its aggressive
stopping rule produces a comparatively concentrated slope-error distribution.
The raw-FPLS mean of $334.298$ does not describe its typical fit: it is driven
by a very small number of catastrophic power-basis solves.

Prediction favours FPCR and Arnoldi FPLS.  Their median MSPEs are $1.112$ and
$1.113$, respectively, only slightly above the irreducible noise floor of
$1$.  Raw FPLS has the same central behaviour, whereas CG--FPLS has median
MSPE $1.230$.  Thus, CG's lower unweighted mean slope error does not imply a
lower covariance-weighted prediction risk: errors in the dominant directions
receive greater weight in MSPE.  This divergence supplies the concentrated-
signal baseline for Model~2, where substantial slope signal extends across
five leading directions.

\subsection{Model 2: a less compressible slope}

Model~2 sets the first five slope coefficients to $4$ without changing the
predictor score variances.  Because the first coefficient already equals $4$
in Model~1, the actual change is an increase in coefficients $2$--$5$.  Every
procedure consequently tends to select a richer model and incurs substantially
larger slope loss.  CG--FPLS has the smallest median and mean ISE, $5.272$ and
$5.506$.  FPCR follows with mean ISE $5.662$, while
Arnoldi FPLS has mean ISE $6.065$.  The differences in their medians are much
smaller: $5.272$, $5.365$, and $5.345$, respectively.  The bias--variance
decomposition below shows that this is the clearest case in which the early
CG rule trades additional bias for reduced slope variance in the ordinary
$L^2$ norm.

The prediction ordering is different.  Arnoldi FPLS has the smallest mean and
median MSPE, $1.442$ and $1.419$, followed closely by the typical raw-FPLS
fit.  FPCR has mean MSPE $1.495$, whereas CG--FPLS reaches $1.999$.  The
directions suppressed by the early CG rule therefore matter materially for
the response even though retaining them increases unweighted slope variance.
Model~2 therefore makes the separation between slope recovery and prediction
especially visible.  Model~3 supplies the complementary intervention: the
slope returns to the baseline while the leading predictor variances change.

\subsection{Model 3: stronger leading predictor directions}

Model~3 keeps the slope of Model~1 and sets the first five score variances to
$2$.  Since the first score variance is already $2$ in Model~1, only score
variances $2$--$5$ are raised.  The resulting stronger information in these
directions is accompanied by a reduction in CG--FPLS mean ISE from
$0.633$ in Model~1 to $0.447$, the best value in Model~3.  FPCR and Arnoldi
FPLS have mean ISEs of $0.615$ and $0.670$.  The median differences are more
modest: $0.426$ for CG--FPLS, approximately $0.446$ for both Krylov
least-squares implementations, and $0.463$ for FPCR.

Prediction again favours the less aggressively stopped procedures.  FPCR has
mean MSPE $1.127$, Arnoldi FPLS $1.131$, and CG--FPLS $1.221$; their medians
are $1.113$, $1.114$, and $1.176$.  Raising the leading score variances
improves identification in the ordinary slope norm, but it also increases the
prediction penalty attached to estimation error in those directions.  The
comparison between Models~1 and 3 therefore shows that a change in the
score-variance sequence can affect ISE and MSPE differently even when the true
slope is identical.  The selected component counts are consequently needed to
show how each tuning rule converts these changes in signal and covariance into
regularisation.

\section{Selected regularisation levels}
\label{sec:results-components}

Table~\ref{tab:results-components} shows how each complete procedure converts
the data into a regularisation level.  The response-informed Krylov methods
generally use far fewer components than FPCR.  This does not mean that one
Krylov component equals one principal component; it reflects the fact that a
Krylov direction combines covariance information with the response moment.

\begin{table}[htbp]
    \centering
    \small
    \caption[Distributions of selected component counts]{Distribution of the selected component count over $5{,}000$
    replications.  The maximum records rare selections and should be read
    together with the numerical diagnostics.}
    \label{tab:results-components}
    \begin{tabular}{@{}llrrrr@{}}
        \toprule
        Model & Method & Mean & Median & $p_{90}$ & Maximum \\
        \midrule
        Model 1 & CG--FPLS       & 1.098 & 1 & 1 & 2 \\
                & Raw FPLS       & 2.134 & 2 & 2 & 57 \\
                & Arnoldi FPLS   & 2.010 & 2 & 2 & 14 \\
                & FPCR           & 5.840 & 5 & 11 & 42 \\
        \addlinespace
        Model 2 & CG--FPLS       & 2.886 & 3 & 3 & 3 \\
                & Raw FPLS       & 7.820 & 5 & 6 & 70 \\
                & Arnoldi FPLS   & 4.808 & 5 & 6 & 16 \\
                & FPCR           & 18.350 & 17 & 28 & 60 \\
        \addlinespace
        Model 3 & CG--FPLS       & 1.833 & 2 & 2 & 3 \\
                & Raw FPLS       & 2.956 & 3 & 3 & 68 \\
                & Arnoldi FPLS   & 2.845 & 3 & 3 & 11 \\
                & FPCR           & 7.057 & 6 & 11 & 46 \\
        \bottomrule
    \end{tabular}
\end{table}

The component counts respond systematically to the controlled changes.
Model~2 exhibits an upward shift consistent with its predictive signal being
less compressible after coefficients $2$--$5$ are increased.  This is not a
one-to-one correspondence between basis coefficients and Krylov components,
because each Krylov direction combines several basis directions.  CG--FPLS
selects three components in $88.6\%$
of its Model~2 replications, compared with one component in $90.2\%$ of
Model~1 replications and two components in $83.2\%$ of Model~3 replications.
FPCR selects many more components---a median of $17$ in Model~2---because its
ordering is determined by predictor variance rather than by the cross-moment
with the response.

All $15{,}000$ CG--FPLS fits cross the plug-in discrepancy threshold before
the upper limit, so none is assigned $m=70$ by the fallback rule.  The internal
moment-GCV pilot used to scale that threshold selects a median of $70$
components in all three models and produces a mean residual variance estimate
of approximately $0.300$, although the data-generating error variance is $1$.
Because the pilot usually fits $70$ components and divides its residual sum of
squares by $n$ without a degrees-of-freedom correction, a value near
$1-70/100=0.30$ is consistent with an overfitting or degrees-of-freedom effect.
The pilot is therefore not an accurate stand-alone noise estimator in this
finite design.  Since the threshold is proportional to $\widehat\sigma$, this
downward bias makes crossing harder and can only delay stopping along the same
residual path.  The observed one-to-three-component indices therefore reflect
rapid residual contraction, but the risk results validate only the implemented
plug-in procedure, not the oracle discrepancy theorem.  The pilot remains part
of the released-code procedure and is not retuned after observing this
discrepancy.

The raw and Arnoldi component distributions coincide in their central region
but diverge sharply in their maxima.  Raw FPLS occasionally selects dimensions
as high as $57$, $70$, and $68$, whereas the Arnoldi maxima are $14$, $16$,
and $11$.  Component counts alone cannot determine whether these extremes
reflect statistical complexity or numerical failure, so the next section
links them to the fixed stability diagnostics.

\section{Complete tails and numerical stability}
\label{sec:results-numerical-stability}

The similarity of medians in Table~\ref{tab:results-primary-risk} must not be
confused with equivalent reliability.  Table~\ref{tab:results-raw-instability}
isolates the raw-FPLS stability diagnostic.  A flagged fit is not automatically
non-finite: all reported losses are finite.  Instead, the flag records that
the condition number of the normalised raw Krylov basis or of its response
design exceeds the fixed threshold.  Because raw FPLS then solves normal
equations in those coordinates, its reconstructed slope is not regarded as
numerically reliable for replay-based function diagnostics.

\begin{table}[htbp]
    \centering
    \small
    \caption[Raw-FPLS instability and upper tails]{Numerical instability and complete upper tails of raw FPLS.
    Stable-subset size is the number of replications available for paired
    function-valued replay.}
    \label{tab:results-raw-instability}
    \begin{tabular}{@{}lrrrrr@{}}
        \toprule
        Model & Flagged fits & Rate & Maximum ISE & Maximum MSPE
              & Stable subset \\
        \midrule
        Model 1 & 20  & 0.40\% & $1.668\times10^{6}$ & $1.005\times10^{5}$ & 4,980 \\
        Model 2 & 409 & 8.18\% & $3.482\times10^{3}$ & $1.578\times10^{2}$ & 4,591 \\
        Model 3 & 16  & 0.32\% & $2.203\times10^{3}$ & $9.737\times10^{1}$ & 4,984 \\
        \bottomrule
    \end{tabular}
\end{table}

Model~2 is most exposed, and this exposure coincides with typical FPLS paths
of approximately five rather than two or three Krylov directions.  Its raw
instability rate is $8.18\%$, compared with $0.40\%$ and $0.32\%$ in Models~1
and 3.  In Model~1, only one raw ISE exceeds one hundred times its
model--method median, but that
single value reaches approximately $1.67$ million.  This is why raw FPLS has a
median ISE of $0.539$ and a mean of $334.298$ in the same model.  Robust
summaries alone would conceal a practically important failure mode, while the
mean alone would conceal the method's ordinarily competitive behaviour.

The combined Arnoldi--QR implementation removes the catastrophic numerical
tail in this experiment, not all statistical tail variation.  Arnoldi FPLS
has no numerical-instability flags,
but its maximum ISEs are $36.35$, $77.41$, and $30.23$ across the three
models.  These values are compatible with ordinary sample and tuning
variability and do not coincide with a loss of Arnoldi-basis orthogonality.
The corresponding FPCR maxima are
$11.53$, $28.23$, and $16.30$, while the CG--FPLS maxima are $1.82$, $16.79$,
and $1.25$.  The especially compact CG tails are consistent with its strong
early-stopping regularisation; the bias--variance decomposition next shows how
much variance reduction this achieves and what bias it incurs.

\section{Bias--variance mechanisms}
\label{sec:results-bias-variance}

Figure~\ref{fig:results-bias-variance-all} decomposes the losses of the three
fully replayable methods over all replications.  For raw FPLS, only the frozen
archive total is displayed in the full-sample panels: a decomposition into a
mean function and dispersion would require reconstructing the $445$
platform-sensitive slopes.  The same restriction applies to the replay-based
basis-coefficient distributions, pointwise distributions, and normal
quantile--quantile plots in Appendix~\ref{app:simulation-diagnostics}.
Table~\ref{tab:results-bias-variance} gives the numerical components for
CG--FPLS, Arnoldi FPLS, and FPCR.  Prediction bias and
variance are weighted by the independent-score variances $\lambda_j$.  This is
the exact excess-prediction-risk geometry of the discrete data-generating
process; the irreducible variance $\sigma^2=1$ must be added to their sum.

\begin{figure}[htbp]
    \centering
    \includegraphics[width=\textwidth]{bias_variance_decomposition.pdf}
    \caption[Full-archive bias--variance decomposition]{Monte Carlo bias--variance decomposition over all $5{,}000$
    replications.  Each row is one model; slope risk is on the left and
    prediction risk on the right.  Bars show squared bias, variance, and irreducible noise on
    logarithmic scales; black diamonds show total risk.  Full raw-FPLS totals
    are retained from the archive, but their components are intentionally not
    reconstructed for numerically unstable fits.}
    \label{fig:results-bias-variance-all}
\end{figure}

\begin{table}[htbp]
    \centering
    \scriptsize
    \caption[Bias--variance components for replayable methods]{Bias--variance components for the fully replayable methods.  The
    prediction columns exclude the irreducible noise variance of $1$.}
    \label{tab:results-bias-variance}
    \begin{tabular}{@{}llrrrr@{}}
        \toprule
        & & \multicolumn{2}{c}{Slope risk}
          & \multicolumn{2}{c}{Excess prediction risk} \\
        \cmidrule(lr){3-4}\cmidrule(l){5-6}
        Model & Method & Squared bias & Variance & Squared bias & Variance \\
        \midrule
        Model 1 & CG--FPLS     & 0.1069 & 0.5259 & 0.1102 & 0.1604 \\
                & Arnoldi FPLS & 0.0082 & 0.7520 & 0.0132 & 0.1163 \\
                & FPCR         & 0.0098 & 0.6602 & 0.0164 & 0.1102 \\
        \addlinespace
        Model 2 & CG--FPLS     & 1.2263 & 4.2796 & 0.4733 & 0.5224 \\
                & Arnoldi FPLS & 0.0990 & 5.9656 & 0.0491 & 0.3903 \\
                & FPCR         & 0.2419 & 5.4197 & 0.1106 & 0.3857 \\
        \addlinespace
        Model 3 & CG--FPLS     & 0.0183 & 0.4282 & 0.0346 & 0.1883 \\
                & Arnoldi FPLS & 0.0055 & 0.6642 & 0.0098 & 0.1230 \\
                & FPCR         & 0.0072 & 0.6078 & 0.0135 & 0.1141 \\
        \bottomrule
    \end{tabular}
\end{table}

The mechanism is consistent in all three models.  CG--FPLS substantially
reduces slope variance but accepts more bias.  In Model~1, for example, its
slope variance is $0.526$, compared with $0.752$ for Arnoldi FPLS and $0.660$
for FPCR, while its squared bias is $0.107$ rather than approximately $0.01$.
This exchange lowers CG--FPLS mean ISE relative to the other replayable
methods.  In prediction geometry, however, the CG squared bias and variance
are both larger: $0.110$ and $0.160$, compared with $0.013$ and $0.116$ for
Arnoldi FPLS.  Adding the common noise variance gives expected MSPEs of
approximately $1.271$ and $1.130$.

Model~2 sharpens the same distinction.  CG--FPLS reduces slope variance by
about $1.69$ relative to Arnoldi FPLS, more than compensating for its larger
unweighted squared bias.  After covariance weighting, CG has prediction
squared bias $0.473$ and variance $0.522$, both above the Arnoldi values of
$0.049$ and $0.390$.  The procedure can therefore be preferable for the
ordinary slope norm and inferior for response prediction without any
contradiction.  Model~3 repeats this pattern on a smaller scale.  Early
stopping is not uniformly ``more biased but less variable'' independently of
the norm: the relevant bias and variance must be evaluated in the geometry of
the target loss.

\section{Raw and Arnoldi FPLS on a paired stable subset}
\label{sec:results-raw-arnoldi}

The preceding decomposition combines differences in statistical
regularisation with differences in numerical implementation.  To isolate the
latter, Figure~\ref{fig:results-bias-variance-stable} and
Table~\ref{tab:results-stable-pair} retain only replications in which the raw
Krylov basis remains stable and apply that same subset to all four methods.
This paired restriction is diagnostic only; it does not replace the primary
full-archive comparison.

\begin{figure}[htbp]
    \centering
    \includegraphics[width=\textwidth]{bias_variance_decomposition_raw_stable_subset.pdf}
    \caption[Bias--variance decomposition on the paired stable subset]{Bias--variance decomposition on the paired subset for which raw
    FPLS satisfies its numerical-stability rule.  All four methods are
    evaluated on exactly the same retained replications within a model.
    Each row is one model, with its retained sample size shown; slope risk
    is on the left and prediction risk on the right.  Bar colours and black
    total-risk diamonds have the same meaning as in
    Figure~\ref{fig:results-bias-variance-all}.}
    \label{fig:results-bias-variance-stable}
\end{figure}

\begin{table}[htbp]
    \centering
    \small
    \caption[Raw and Arnoldi errors on the paired stable subset]{Mean errors of raw and Arnoldi FPLS on paired numerically stable
    raw-FPLS replications.}
    \label{tab:results-stable-pair}
    \begin{tabular}{@{}lrrrrr@{}}
        \toprule
        & & \multicolumn{2}{c}{Mean slope ISE}
          & \multicolumn{2}{c}{Mean test MSPE} \\
        \cmidrule(lr){3-4}\cmidrule(l){5-6}
        Model & $R_{\mathrm{stable}}$ & Raw & Arnoldi & Raw & Arnoldi \\
        \midrule
        Model 1 & 4,980 & 0.7085 & 0.7166 & 1.1279 & 1.1282 \\
        Model 2 & 4,591 & 5.5739 & 5.6121 & 1.4271 & 1.4282 \\
        Model 3 & 4,984 & 0.6384 & 0.6384 & 1.1302 & 1.1302 \\
        \bottomrule
    \end{tabular}
\end{table}

On the stable subsets, raw and Arnoldi FPLS are nearly indistinguishable.
Their Model~3 mean ISEs differ by less than $3\times10^{-7}$ and their mean
MSPEs by less than $10^{-8}$.  Small differences in Models~1 and 2 arise
because each implementation minimises its own finite-precision
cross-validation curve, so a small number of selected dimensions differ even
before the instability threshold is reached.  The comparison supports the
equivalence boundary established in Chapter~\ref{chap:methods}: when the raw
Krylov basis remains usable, it and the Arnoldi basis approximate the same
response-least-squares solution; the practical advantage of Arnoldi is that
this equivalence remains numerically accessible farther along the path when
the orthogonal basis is paired with the response-space QR solve.

Stabilisation therefore extends how far the shared FPLS target remains
numerically accessible; it does not create a different statistical target or
guarantee a smaller loss in every sample.  CGS2 supplies reliable coordinates
for the same nested spaces, while QR avoids response-space normal equations.
The coefficientwise and pointwise distributions considered next show what
each tuning rule then retains or suppresses in the recovered slope.

\section{Distribution of the recovered slope}
\label{sec:results-slope-distributions}

The global error measures are complemented by coefficientwise and pointwise
diagnostics in Appendix~\ref{app:simulation-diagnostics}.  Figures
\ref{fig:app-coefficient-distributions-model1}--
\ref{fig:app-coefficient-distributions-model3} show the first three basis
coefficients, and Figures~\ref{fig:app-pointwise-distributions-model1}--
\ref{fig:app-pointwise-distributions-model3} show the estimated slope at five
locations.  The raw-FPLS panels use only archived fits satisfying the raw
stability rule; the other methods use all $5{,}000$ replications.

These coefficients are recovered from each slope vector on the finite grid,
rather than computed as continuum inner products.  If
$\boldsymbol\Phi\in\mathbb{R}^{J\times T}$ has entries
$\Phi_{jt}=\phi_j(s_t)$, the plotted coefficient vector is
\begin{equation}
    \widehat{\boldsymbol b}
    = (\boldsymbol\Phi\boldsymbol\Phi^{\mathsf T})^{-1}
      \boldsymbol\Phi\widehat{\boldsymbol\beta},
    \label{eq:discrete-basis-coefficient-recovery}
\end{equation}
implemented by a Moore--Penrose pseudoinverse.  On the endpoint-inclusive
grid this least-squares recovery is not identical to simply taking discrete
inner products with the nominal basis functions.

The coefficient distributions make the effect of early stopping concrete.
In Model~1, CG--FPLS estimates the second coefficient with mean $0.366$ when
its true value is $0.616$, and the third with mean $0.071$ rather than $0.206$.
Its first coefficient is also shrunk, but much less severely: $3.867$ rather
than $4$.  Since $90.2\%$ of CG fits stop at one component, the procedure
captures the dominant supervised direction and deliberately leaves much of
the weaker structure unresolved.  Raw FPLS, Arnoldi FPLS, and FPCR have
smaller bias in these later coefficients but wider functional variation.

The pattern changes with the design.  In Model~2, where all first five slope
coefficients equal $4$, CG--FPLS means for the first three are $3.930$,
$3.848$, and $3.887$.  In Model~3, the larger leading score variances are
associated with CG--FPLS means of $0.599$ and $0.198$ for coefficients two and
three, close to the true values $0.616$ and $0.206$.  This is consistent with
the marked improvement in CG slope ISE from Model~1 to Model~3.

The pointwise distributions tell the same story after the coefficients have
been aggregated into functions.  CG--FPLS is generally narrower but can be
visibly shifted from the truth, particularly in Model~1.  Model~2 has the
widest distributions for every method, a pattern consistent with the less
compressible, oscillatory slope and the use of several estimated directions.
The normal quantile--quantile plots in
Figures~\ref{fig:app-coefficient-qq-model1}--
\ref{fig:app-coefficient-qq-model3} are approximately linear through their
centres for most method--coefficient combinations.  Departures in the tails
and bends for some CG coordinates reflect the mixture induced by a discrete,
data-dependent stopping rule; they should not be interpreted as evidence
that the complete functional estimator has a simple Gaussian distribution.

Across the loss, complexity, stability, and function-level evidence, two
distinctions govern the comparison.  Aggressive CG stopping stabilises
ordinary slope estimation but can leave covariance-weighted signal unresolved;
Arnoldi does not change the full-rank APLS target, but keeps that target
numerically accessible when raw powers become nearly collinear.  Any ranking
therefore depends both on the statistical target and on whether the computation
reliably realises the intended estimator.

These conclusions are conditional on dense error-free curves, Gaussian
scores, $n=100$, the three chosen designs, and the stated tuning rules.  They
identify mechanisms rather than a universal ranking.  The compact
spectroscopy application in Chapter~\ref{chap:spectroscopy-application}
provides a first assessment beyond this controlled setting, while
Appendix~\ref{app:spectroscopy-supplement} retains the complete empirical
diagnostics and sensitivity evidence.

That dependence changes again for inference.  Estimation can benefit from the
bias introduced by a small $m$, whereas
Chapter~\ref{chap:inference-methodology} requires the CG residual to be
negligible relative to $n^{-1/2}$ sampling uncertainty.  The early stopping
used here therefore cannot be transferred unchanged to hypothesis testing.
